\documentclass[10pt,a4paper]{article}
\usepackage[margin=2.2cm]{geometry}
\usepackage[T1]{fontenc}
\usepackage{lmodern}
\usepackage{microtype}
\usepackage{amsmath,amssymb}
\usepackage{booktabs}
\usepackage{enumitem}
\usepackage[font=small,labelfont=bf]{caption}
\usepackage[colorlinks=true,linkcolor=blue!50!black,urlcolor=blue!50!black]{hyperref}

\setlength{\parindent}{0pt}
\setlength{\parskip}{0.4em}
\setlist{nosep,leftmargin=*,topsep=3pt}
\renewcommand{\arraystretch}{1.12}

\newcommand{\nbar}{\bar n}
\newcommand{\dtil}{\tilde\delta}
\newcommand{\B}[1]{\textbf{#1}}

\title{\textbf{Point-process TDA: status and results on the new dataset (DV3)}}
\date{\today}

\begin{document}
\maketitle

% ============================================================================
\section{Where things stand}

\begin{itemize}
  \item A new dataset (DV3) was generated from scratch so that performance can be reported \emph{per regime} (how far from Poisson, how many points) instead of as one averaged number.
  \item \textbf{Done:} the classical baselines (a neural network reading the standard summary functions), for parameter estimation and 5-way model classification; and the first persistence-homology (PH) feature, the $H_0$ persistence image of the DTM filtration ($k=5$). Every number is an average over 10 training seeds.
  \item \textbf{Main result so far:} the classical baselines win everywhere. For parameter estimation, the single best model on every family is classical. For classification, the best classical model (all four summary functions) beats the PH model by 7 points of accuracy (0.711 vs.\ 0.641).
  \item \textbf{Not yet run:} the PH $H_1$ persistence image, Betti curves and persistence landscapes; DTM $k=10, 15$ and the Rips filtration; and the classical non-learned references (minimum-contrast fitting, envelope tests). The $H_1$ image was tried, but it currently fails to train for a technical reason (its pixel values are $\sim$10$\times$ smaller than $H_0$'s and get ignored by the network); it will be rerun with input normalisation.
\end{itemize}

% ============================================================================
\section{Data}

\textbf{Models.} Five stationary models on the unit square: Poisson (complete spatial randomness, CSR), Thomas (clustered), nested Thomas (clusters of clusters), Mat\'ern type~II (hard-core, i.e.\ repulsive) and the log-Gaussian Cox process (LGCP). The expected number of points $\nbar$ ranges from 100 to 800.

\textbf{How the distance from Poisson is measured ($\dtil$).}
\begin{itemize}
  \item The classical way to test a pattern against Poisson is to compare its $L$-function ($L(r)-r$ is zero for Poisson) with its Poisson expectation, taking the largest standardised deviation over a range of radii (a global envelope test).
  \item For every parameter setting we compute that largest deviation for the model's \emph{theoretical} $L$-function, and divide it by the test's 5\% critical value at the same $\nbar$. That ratio is $\dtil$.
  \item So $\dtil = 0$ is exactly Poisson, $\dtil \approx 1$ is the edge of what the classical test can detect at the 5\% level, and larger values are more obviously non-Poisson. At a fixed shape, $\dtil$ grows roughly like $\sqrt{\nbar}$.
  \item Results are broken down into $\dtil$ bands: $[0,0.5)$ and $[0.5,1)$ are "below the classical detection threshold", $\ge 2$ is "clearly structured".
\end{itemize}

\textbf{How the sets are built.}
\begin{itemize}
  \item \textbf{Training:} 10{,}000 patterns per model, parameters drawn at random from a prior (the prior deliberately puts roughly a quarter to a half of its mass at $\dtil \le 1$, depending on the model). A fixed 85/15 train/validation split, identical for every method and seed.
  \item \textbf{Test set A (average performance):} 5{,}000 fresh patterns per model from the same prior. This is the headline number.
  \item \textbf{Test set B (fixed settings):} a grid of fixed parameter settings (combinations of $\nbar$, cluster size and $\dtil$), with 400 replicate patterns each. It measures error at specific, controlled settings.
  \item \textbf{Test set C (paths to Poisson):} for each model, the shape is held fixed while the strength of the structure is turned down in 16 steps until the model is exactly Poisson (300 replicates per step, at two values of $\nbar$). Exact Poisson patterns are added and used to calibrate the 5\% detection threshold. It measures how performance degrades as the pattern approaches Poisson.
\end{itemize}

% ============================================================================
\section{Methods compared}

\begin{itemize}
  \item \textbf{Classical baselines:} the 1-D CNN of Vihrs (2022), reading summary functions sampled on a radius grid, plus the point count $n(x)$. Four input sets: $L$ alone (the published baseline), $L+F+G+J$ (all four standard summary functions), $G$ alone, $F$ alone. $F$, $G$, $J$ were tried on two radius grids ($r_{\max}=0.25$ and $0.08$); the grid was chosen on the validation set, never on the test sets.
  \item \textbf{PH:} a CNN on the persistence image of the $H_0$ diagram of the DTM filtration ($k=5$), plus $n(x)$ for parameter estimation. It is one filtration and one homology dimension, with no fusion and no summary functions mixed in.
  \item \textbf{Training:} the same for all: early stopping on validation loss, 10 seeds (9371--9380). Tables report mean $\pm$ standard deviation over seeds.
\end{itemize}

% ============================================================================
\section{Parameter estimation (inference)}

\begin{itemize}
  \item \textbf{Metric:} normalised loss = mean squared error of the log-parameters, standardised with the training statistics, averaged over the model's parameters. Lower is better; $\approx 1$ means predicting the training mean. Values are \emph{not} comparable across models.
  \item \textbf{Parameters estimated:} Thomas $(\kappa,\mu,\sigma)$; nested $(\kappa,\mu_1,\mu_2,\sigma_1,\sigma)$; Mat\'ern II $(\lambda_p,R)$; LGCP $(\nbar,\sigma^2,s)$.
  \item \textbf{Bold} marks the single best model per family on test set A.
\end{itemize}

\begin{table}[ht]
\centering\small
\begin{tabular}{@{}ll ccc@{}}
\toprule
Family & Model & A (prior) & B (fixed settings) & C (paths to Poisson)\\
\midrule
Thomas & $L$ & 0.1365 $\pm$ 0.0007 & 0.1371 $\pm$ 0.0054 & 0.2818 $\pm$ 0.0208\\
 & \B{$L{+}F{+}G{+}J$} & \B{0.1364 $\pm$ 0.0015} & \B{0.1406 $\pm$ 0.0056} & \B{0.2749 $\pm$ 0.0140}\\
 & $G$ & 0.3308 $\pm$ 0.0031 & 0.4021 $\pm$ 0.0134 & 0.6438 $\pm$ 0.0442\\
 & $F$ & 0.1926 $\pm$ 0.0020 & 0.2033 $\pm$ 0.0079 & 0.4096 $\pm$ 0.0401\\
 & PH image $H_0$ & 0.1771 $\pm$ 0.0045 & 0.1931 $\pm$ 0.0078 & 0.3923 $\pm$ 0.0319\\
\midrule
Nested & $L$ & 0.3645 $\pm$ 0.0023 & 0.2045 $\pm$ 0.0102 & 0.1689 $\pm$ 0.0128\\
 & \B{$L{+}F{+}G{+}J$} & \B{0.3582 $\pm$ 0.0029} & \B{0.1951 $\pm$ 0.0205} & \B{0.1713 $\pm$ 0.0237}\\
 & $G$ & 0.4363 $\pm$ 0.0015 & 0.1960 $\pm$ 0.0141 & 0.1572 $\pm$ 0.0163\\
 & $F$ & 0.4683 $\pm$ 0.0049 & 0.2643 $\pm$ 0.0143 & 0.2544 $\pm$ 0.0242\\
 & PH image $H_0$ & 0.4540 $\pm$ 0.0111 & 0.2531 $\pm$ 0.0234 & 0.2515 $\pm$ 0.0404\\
\midrule
Mat\'ern II & $L$ & 0.0184 $\pm$ 0.0008 & 0.0166 $\pm$ 0.0008 & 0.0178 $\pm$ 0.0008\\
 & \B{$L{+}F{+}G{+}J$} & \B{0.0180 $\pm$ 0.0003} & \B{0.0163 $\pm$ 0.0005} & \B{0.0172 $\pm$ 0.0007}\\
 & $G$ & 0.0185 $\pm$ 0.0003 & 0.0160 $\pm$ 0.0003 & 0.0174 $\pm$ 0.0004\\
 & $F$ & 0.0995 $\pm$ 0.0044 & 0.0922 $\pm$ 0.0042 & 0.0934 $\pm$ 0.0050\\
 & PH image $H_0$\,$^\dagger$ & 0.1002 $\pm$ 0.0658 & 0.0868 $\pm$ 0.0588 & 0.0930 $\pm$ 0.0647\\
\midrule
LGCP & \B{$L$} & \B{0.2203 $\pm$ 0.0011} & \B{0.1889 $\pm$ 0.0042} & \B{0.1051 $\pm$ 0.0053}\\
 & $L{+}F{+}G{+}J$ & 0.2217 $\pm$ 0.0027 & 0.1902 $\pm$ 0.0054 & 0.1052 $\pm$ 0.0071\\
 & $G$ & 0.3084 $\pm$ 0.0013 & 0.2637 $\pm$ 0.0051 & 0.1055 $\pm$ 0.0031\\
 & $F$ & 0.2775 $\pm$ 0.0012 & 0.2357 $\pm$ 0.0056 & 0.0975 $\pm$ 0.0045\\
 & PH image $H_0$\,$^\dagger$ & 0.3465 $\pm$ 0.1597 & 0.2611 $\pm$ 0.0874 & 0.1422 $\pm$ 0.0665\\
\bottomrule
\end{tabular}
\caption{Parameter estimation: normalised loss (lower is better), mean $\pm$ SD over 10 seeds. $^\dagger$\,Unstable across seeds; see the notes below.}
\label{tab:params}
\end{table}

\begin{table}[ht]
\centering\small
\begin{tabular}{@{}ll cccccc@{}}
\toprule
 & & \multicolumn{6}{c}{Test set A, by distance from Poisson $\dtil$}\\
\cmidrule(l){3-8}
Family & Model & [0, 0.5) & [0.5, 1) & [1, 2) & [2, 4) & [4, 8) & $\ge 8$\\
\midrule
Thomas & $L$ & 0.290 & 0.268 & 0.189 & 0.072 & 0.033 & 0.021\\
 & \B{$L{+}F{+}G{+}J$} & \B{0.283} & \B{0.278} & \B{0.192} & \B{0.070} & \B{0.033} & \B{0.020}\\
 & PH image $H_0$ & 0.338 & 0.338 & 0.274 & 0.134 & 0.051 & 0.022\\
\midrule
Nested & $L$ & 0.503 & 0.528 & 0.465 & 0.309 & 0.253 & 0.179\\
 & \B{$L{+}F{+}G{+}J$} & \B{0.504} & \B{0.535} & \B{0.457} & \B{0.305} & \B{0.226} & \B{0.171}\\
 & PH image $H_0$ & 0.597 & 0.593 & 0.548 & 0.441 & 0.359 & 0.239\\
\midrule
Mat\'ern II & $L$ & 0.033 & 0.017 & 0.013 & 0.009 & 0.010 & 0.010\\
 & \B{$L{+}F{+}G{+}J$} & \B{0.032} & \B{0.016} & \B{0.013} & \B{0.010} & \B{0.009} & \B{0.007}\\
 & PH image $H_0$ & 0.149 & 0.047 & 0.079 & 0.110 & 0.058 & 0.073\\
\midrule
LGCP & \B{$L$} & \B{0.275} & \B{0.290} & \B{0.252} & \B{0.180} & \B{0.142} & \B{0.139}\\
 & $L{+}F{+}G{+}J$ & 0.281 & 0.289 & 0.256 & 0.179 & 0.141 & 0.138\\
 & PH image $H_0$ & 0.438 & 0.338 & 0.297 & 0.270 & 0.299 & 0.388\\
\bottomrule
\end{tabular}
\caption{Parameter estimation on test set A by $\dtil$ band (normalised loss, mean over 10 seeds). Bold = the family's best model from Table~\ref{tab:params}.}
\label{tab:params-delta}
\end{table}

\begin{itemize}
  \item \textbf{A classical model wins on every family.} $L{+}F{+}G{+}J$ is best on Thomas, nested and Mat\'ern II; $L$ alone is best on LGCP. On Thomas, Mat\'ern II and LGCP, $L$ and $L{+}F{+}G{+}J$ are within about one standard deviation of each other, so adding $F,G,J$ to $L$ buys little; only on nested is the gain clear (0.358 vs.\ 0.365).
  \item \textbf{PH image $H_0$ is behind the best classical model on every family:} Thomas 0.177 vs.\ 0.136, nested 0.454 vs.\ 0.358. It is better than $G$ alone on Thomas and than $F$ alone on Thomas and nested.
  \item \textbf{Where PH loses:} near Poisson. On Thomas the gap to the best model is 0.05--0.08 for $\dtil < 4$, and it disappears for strongly clustered patterns ($\dtil \ge 8$: 0.022 vs.\ 0.020).
  \item $^\dagger$\,\textbf{PH is unstable on Mat\'ern II and LGCP.} On Mat\'ern II, 6 seeds reach 0.046--0.054 but 4 stopped early at 0.14--0.19. On LGCP, 7 seeds reach $\approx$0.25 but 3 got stuck at 0.58, the value obtained from $n(x)$ alone. Even the good seeds are well behind the classical models (Mat\'ern II 0.018, LGCP 0.220), so the conclusion does not depend on the failed seeds.
  \item Loss broadly falls as $\dtil$ grows (not for PH on Mat\'ern II and LGCP, whose averages mix good and failed seeds). On Mat\'ern II the classical models are almost exact ($\approx$0.01--0.03).
\end{itemize}

% ============================================================================
\section{Classification}

\begin{itemize}
  \item \textbf{Task:} 5-way classification (Poisson, Thomas, nested, Mat\'ern II, LGCP). All models are trained on the same pooled training set (50{,}000 patterns).
  \item \textbf{Metric:} accuracy (higher is better; chance is 0.20). Test set A has 5{,}000 patterns per model, so its accuracy is the headline. B and C mix the models unevenly by design and are shown for completeness.
  \item \textbf{Bold} marks, for each family, the single model that recognises it best (recall on A).
\end{itemize}

\begin{table}[ht]
\centering\small
\begin{tabular}{@{}l ccc ccccc@{}}
\toprule
 & \multicolumn{3}{c}{Overall accuracy} & \multicolumn{5}{c}{Test set A: recall per family}\\
\cmidrule(lr){2-4}\cmidrule(l){5-9}
Model & A & B & C & Poisson & Thomas & Nested & Mat\'ern II & LGCP\\
\midrule
$L$ & 0.686 $\pm$ 0.002 & 0.617 & 0.614 & 0.813 & 0.503 & 0.599 & \B{0.906} & 0.610\\
$L{+}F{+}G{+}J$ & 0.711 $\pm$ 0.002 & 0.659 & 0.635 & \B{0.836} & \B{0.522} & \B{0.631} & 0.892 & \B{0.673}\\
$G$ & 0.606 $\pm$ 0.003 & 0.504 & 0.479 & 0.806 & 0.425 & 0.443 & 0.855 & 0.502\\
$F$ & 0.573 $\pm$ 0.002 & 0.498 & 0.497 & 0.800 & 0.398 & 0.492 & 0.616 & 0.558\\
PH image $H_0$ & 0.641 $\pm$ 0.005 & 0.585 & 0.571 & 0.808 & 0.445 & 0.625 & 0.666 & 0.663\\
\bottomrule
\end{tabular}
\caption{Classification accuracy, mean over 10 seeds (SD shown for A; B and C SDs are $\le 0.02$).}
\label{tab:classify}
\end{table}

\begin{table}[ht]
\centering\footnotesize\setlength{\tabcolsep}{3.5pt}
\begin{tabular}{@{}l cccccc c c cccccc@{}}
\toprule
 & \multicolumn{6}{c}{Accuracy on A, by $\dtil$ (structured models)} & & & \multicolumn{6}{c}{Detection power on C, by $\dtil$}\\
\cmidrule(lr){2-7}\cmidrule(l){10-15}
Model & [0,.5) & [.5,1) & [1,2) & [2,4) & [4,8) & $\ge$8 & & size & [0,.5) & [.5,1) & [1,2) & [2,4) & [4,8) & $\ge$8\\
\midrule
$L$ & 0.349 & 0.505 & 0.676 & 0.790 & 0.845 & 0.880 & & 0.049 & 0.274 & 0.581 & 0.899 & 0.990 & 0.999 & 1.000\\
$L{+}F{+}G{+}J$ & 0.327 & 0.496 & 0.689 & 0.841 & 0.917 & 0.945 & & 0.049 & 0.259 & 0.568 & 0.884 & 0.990 & 0.999 & 1.000\\
$G$ & 0.263 & 0.380 & 0.538 & 0.689 & 0.793 & 0.784 & & 0.051 & 0.171 & 0.345 & 0.652 & 0.898 & 0.988 & 1.000\\
$F$ & 0.146 & 0.303 & 0.509 & 0.682 & 0.792 & 0.801 & & 0.048 & 0.133 & 0.380 & 0.647 & 0.935 & 0.999 & 1.000\\
PH image $H_0$ & 0.151 & 0.406 & 0.632 & 0.791 & 0.869 & 0.927 & & 0.051 & 0.155 & 0.445 & 0.766 & 0.979 & 0.999 & 1.000\\
\bottomrule
\end{tabular}
\caption{Left: classification accuracy on A by $\dtil$ band, pooling the four non-Poisson models. Right: the classifiers reused as tests of "Poisson vs.\ not Poisson" (score $=1-P(\text{Poisson})$), with the threshold set at 5\% on exact-Poisson patterns; "size" is the achieved false-alarm rate on held-out Poisson patterns, then the power on the paths to Poisson.}
\label{tab:classify-delta}
\end{table}

\begin{itemize}
  \item \textbf{$L{+}F{+}G{+}J$ is the best classifier:} 0.711 on A, and it recognises four of the five families best. $L$ alone recognises Mat\'ern II best (0.906).
  \item \textbf{PH image $H_0$ (0.641) is third of the five models on A,} 7 points behind $L{+}F{+}G{+}J$ and 4.5 points behind $L$ alone. It is ahead of $G$ and $F$ alone. It is nearly as good as the best on nested (0.625 vs.\ 0.631) and LGCP (0.663 vs.\ 0.673), but clearly worse on Thomas (0.445) and especially Mat\'ern II (0.666 vs.\ 0.89--0.91).
  \item \textbf{By distance from Poisson:} every model is weak below $\dtil = 0.5$ (0.15--0.35, chance being 0.20) and improves steadily. PH is weakest near Poisson (0.151 below $\dtil=0.5$, vs.\ 0.349 for $L$). For strongly structured patterns it overtakes $L$ alone ($\dtil \ge 4$: 0.869--0.927 vs.\ 0.845--0.880), but stays behind $L{+}F{+}G{+}J$.
  \item \textbf{Detecting departures from Poisson:} all models hold the 5\% false-alarm rate (0.048--0.051). $L$ and $L{+}F{+}G{+}J$ are the most powerful where it matters ($\dtil < 2$): about 0.57--0.58 power at $\dtil\in[0.5,1)$, vs.\ 0.445 for PH and 0.35--0.38 for $G$ or $F$ alone. Above $\dtil \approx 4$ every model detects essentially everything.
\end{itemize}

% ============================================================================
\section{Next steps}

\begin{itemize}
  \item Fix the $H_1$ persistence image (normalise the input intensities) and rerun it.
  \item Run the remaining PH features: Betti curves and persistence landscapes, $H_0$ and $H_1$ separately; then DTM $k=10, 15$ and the Rips filtration (these diagrams still have to be computed).
  \item Investigate the unstable PH training on Mat\'ern II and LGCP (seeds that stop early or collapse to the $n(x)$-only solution).
  \item Add the non-learned classical references: minimum-contrast parameter fits and the classical envelope tests for detection.
\end{itemize}

\end{document}
